% Registered A.51: full independent root and DE proof/source/code audits PASS.
\subsection{Quaternion multiplication completes the Heisenberg fixed groups}

\begin{statement}{A.51}{Proved}{Full fixed inertia in the genus-ten Heisenberg cover}
For the curve \(X=C\setminus S\) of A.50,
every coefficient prime, every selected cusp packet,
and every indicated open derived subgroup, the full
geometric pro-prime fixed subgroup is the selected
cusp inertia. Together with A.50, both assertions
of Task A.5 hold for this degree-twenty-seven cover.

More precisely, the same surface \(B\) in
(A.50.2) has rational quaternion multiplication by
\[
                         (-3,2)_{\mathbf Q}.
                                                               \tag{A.51.1}
\]
The Jacobian \(J_C\) has potentially good
reduction at three with special three-rank zero,
and every difference of two of the nine cusp
points is annihilated by nine in \(J_C\).
\end{statement}

\begin{proof}
\paragraph{Identify the polarized surface by an exact cellular calculation.}
We compute the rational cup form, not a numerical
period approximation. Write
\[
 H_3=\{(a,b,c):a,b,c\in\mathbf F_3\},\qquad
 (a,b,c)(a',b',c')=(a+a',b+b',c+c'+ab'),
\]
and put \(x=(1,0,0)\), \(y=(0,1,0)\),
\(c=(0,0,1)=[x,y]\). The Heisenberg quotient
of the rank-two free group is its class-two,
exponent-three quotient. In particular every
epimorphism onto \(H_3\) from two handle
generators has the same kernel up to the evident
deck labelling. The topological cover underlying
(A.50.1), whose degree and deck group were proved
there, is therefore the following square cover
of the square torus.

Label the \(27\) oriented squares by \(g\in H_3\).
The right neighbour is \(gx\), and the upper
neighbour is \(gy\). There are nine original
vertices, labelled by \(g\bmod\langle c\rangle\).
Write \(v(g)\) for such a vertex. The oriented
horizontal and vertical edges are
\[
 h_g:v(g)\longrightarrow v(gx),\qquad
 k_g:v(g)\longrightarrow v(gy).
\]
The oriented boundary of square \(g\) is
\[
                       h_g+k_{gx}-h_{gy}-k_g.       \tag{A.51.2}
\]
Subdivide each square by adding its centre \(m_g\)
and the four radial edges \(r_{g,j}\), directed
from the centre to corner \(j\), for corners
\[
        v(g),\quad v(gx),\quad v(gxy),\quad v(gy)
                     \qquad(j=0,1,2,3).
\]
There are \(36\) vertices, \(162\) edges, and
\(108\) positively oriented triangles. The
boundary of triangle \((g,j)\) is
\[
                   r_{g,j}+e_{g,j}-r_{g,j+1},
      \qquad(e_{g,0},e_{g,1},e_{g,2},e_{g,3})
                   =(h_g,k_{gx},-h_{gy},-k_g),       \tag{A.51.3}
\]
with radial subscripts read modulo four.
These formulas specify every entry of the two
integral cellular boundary matrices.

For clarity the cup signs are also explicit.
Order the nine original vertices lexicographically
before all the centres. For two rational cellular
one-cocycles \(a,b\), a triangle with oriented
outside edge \(e=\epsilon e_+\) contributes
\[
                -\epsilon\,a(e_+)\,b(r_{\rm high})
                                                               \tag{A.51.4}
\]
to \(\langle a\smile b,[C]\rangle\), where
\(r_{\rm high}\) is the radial edge ending at
the larger of the two original vertices. This
is the ordered-triangle cup formula: if the
outside vertices are increasingly ordered, the
ordered triangle is their pair followed by its
centre; reversing their order changes both signs
as in (A.51.4). It also follows on barycentrically
subdividing this ordered triangular complex.
Thus (A.51.4) computes the actual topological
polarization on cohomology.

Deck transformations act by left multiplication
on all labels. A lift of quarter rotation has
\[
 \alpha(a,b,c)=(-b,a,c-ab),\qquad
 h_g\longmapsto k_{\alpha(g)x},\quad
 k_g\longmapsto-h_{\alpha(g)},\quad
 r_{g,j}\longmapsto r_{\alpha(g),j+1}.              \tag{A.51.5}
\]
The formulas respect (A.51.2)--(A.51.3), have
order four, and conjugate the deck generators
by \(x\mapsto y,y\mapsto x^{-1}\).
They give a holomorphic lift of quarter rotation
of the square elliptic curve. It and the lift
(A.50.7) are conjugate by a deck transformation:
both are complements of the odd group \(H_3\)
in the extension by \(C_4\), and coprime
complement conjugacy applies. This comparison is
also why the following projector selects the
same factor \(B\) as A.50. Conjugating the lift
merely replaces \(q_0\) by one of its deck
translates, whose entire span is unchanged.

Let \(U,V,Z,R\) denote the resulting rational
operators on \(H^1(C,\mathbf Q)\), using inverse
pullback so that they form a covariant group
representation. They satisfy
\[
 UV=ZVU,\quad RUR^{-1}=V,\quad RVR^{-1}=U^{-1}.
\]
Put
\[
\begin{split}
 P&=1-\frac{1+Z+Z^2}{3},\qquad
 P_0=\frac{1+U+U^2}{3},\\
 A&=\frac{1+2Z}{3}
       \sum_{j,k=0}^2 Z^{-jk}V^jP_0V^{-k}\,P,
       \qquad T=A^3RP,\\
 P_-&=\frac{P-T}{2},\qquad e=P_0P_- .
\end{split}                                                   \tag{A.51.6}
\]
The matrix units \(V^jP_0V^{-k}\) show that
\(A\) is exactly the Fourier intertwiner
(A.50.11) on the nontrivial-central summand.
In particular
\[
 A^4=P,\quad T^2=P,\quad
 \operatorname{rank}P=18,\quad
 \operatorname{rank}\tfrac{P+T}{2}=6,\quad
 \operatorname{rank}P_-=12,\quad\operatorname{rank}e=4.
                                                               \tag{A.51.7}
\]
The plus part is the span of the actual elliptic
quotient maps from A.50; the minus part is
\(H^1(B^3,\mathbf Q)\). Taking \(P_0\), the
invariants of the noncentral order-three subgroup,
selects one copy of \(H^1(B,\mathbf Q)\).
These identifications use actual rational
endomorphisms of the Jacobian. They do not assign
a new geometric meaning to an arbitrary matrix
subspace.

Here is the complete small output needed from the
exact certificate. On \(eH^1(C,\mathbf Q)\),
there are vectors \(v_1,v_2\) forming a basis
over \(F=\mathbf Q(\zeta_3)\), with \(Z\)
acting as \(\zeta_3\), such that the cup matrix
in the rational basis \((v_1,Zv_1,v_2,Zv_2)\) is
\[
 \Psi_B=
 \begin{pmatrix}
 0&-1/6&0&0\\1/6&0&0&0\\
 0&0&0&1/3\\0&0&-1/3&0
 \end{pmatrix}.                                  \tag{A.51.8}
\]
The file \texttt{a\_heisenberg\_polarization.json}
retains all four cellular cocycles as exact rational
values on the labelled edges. The accompanying
program constructs (A.51.2)--(A.51.4), fixes a
spanning-tree gauge for coboundaries, and computes
a basis of the resulting twenty-dimensional
cocycle space. Its cup determinant is one.
It checks the complete deck and rotation relations,
symplectic preservation, every projector identity
and rank in (A.51.7), and (A.51.8) on the four
retained cocycles. This is a finite exact rational
certificate for the displayed pairing, with all
entries specified by the cellular formulas; it
uses no bounded-degree extrapolation.

Write \(\kappa=Z-Z^2=\sqrt{-3}\). The
corresponding Hermitian form is
\[
 h(v,w)=\frac{\Psi_B(v,\kappa w)}2
                  -\frac{\kappa\Psi_B(v,w)}2.
                                                               \tag{A.51.9}
\]
It is linear in its first variable and conjugate
linear in its second. Formula (A.51.8) gives
\[
                  h=\operatorname{diag}(-1/6,1/3),
                    \qquad-\det h=1/18.             \tag{A.51.10}
\]
Scaling the displayed polarization by a rational
number, or changing the \(F\)-basis, changes this
discriminant only by
an \(F\)-norm, so this particular exact choice
is sufficient for the quaternion calculation.

\paragraph{Construct the rational Hodge endomorphism.}
We also verify the Hodge input explicitly.
For the central automorphism \(z\mapsto\zeta_3z\)
in (A.50.1), the holomorphic eigendifferentials
of nontrivial character have dimensions three
and six. One may use the following direct bases:
\[
\begin{array}{c|l}
\zeta_3^2&dx/z,\quad dx/(yz),\quad x\,dx/(yz)\\
\zeta_3&dx/z^2,\quad x\,dx/z^2,
      \quad x^jdx/(yz^2)\ (0\le j\le3).
\end{array}                                                     \tag{A.51.11}
\]
At a finite branch point, \(x-a\) has order
three and \(z\) order one. At the single point
above infinity, \(x,y,z,dx\) have orders
\(-6,-9,-8,-7\), respectively. These facts
check holomorphicity; the eigenvalues and distinct
functions check independence. Together with the
one invariant elliptic differential they give
all ten differentials. Thus the threefold factor
\(B_3\) in A.50 has \(F\)-signature \((2,1)\),
up to exchanging the two embeddings. The elliptic
quotient differential is
\(q_0^*(dz/s)=4x\,dx/z^2\), of the larger
character. Removing its three deck translates
leaves the surface \(B\) with signature \((1,1)\).
Inverse pullback exchanges the two embeddings
and leaves this balanced signature unchanged.

The Hodge circle on \(H^1(B,\mathbf Q)\)
preserves \(h\), and its determinant over
\(F\) is one: at either embedding its two
circle characters are \(z\) and \(z^{-1}\).
It therefore lies in \(\operatorname{SU}(h)\).
This is the dimension-two Weil-type condition;
no equality with the full special unitary group
is required.

Define the following \(F\)-antilinear rational
operator on the basis of (A.51.10):
\[
                  j(v_1)=-v_2,\qquad j(v_2)=-2v_1,
                      \qquad j(av)=\bar a\,j(v).
                                                               \tag{A.51.12}
\]
A direct two-by-two calculation shows that it
commutes with \(\operatorname{SU}(h)\).
Indeed for a diagonal form \(\operatorname{diag}(a,b)\),
the antilinear operator with matrix
\(\bigl(\begin{smallmatrix}0&-b\\a&0\end{smallmatrix}\bigr)\)
commutes with every determinant-one isometry;
(A.51.12) is six times this operator for
\(a=-1/6,b=1/3\).
It follows that \(j\) is a rational Hodge
endomorphism. Weight-one Hodge realization is
fully faithful for complex abelian varieties,
by \cite[Theorems 4.4 and 7.12, pp.~489 and 508]{fghMilneSVModuli}.
After clearing lattice denominators this gives
an actual algebraic endomorphism of the complex
abelian variety. Endomorphisms of a number-field
abelian variety descend to a finite algebraic
extension; enlarge the field to define these
finitely many endomorphisms. If one works
contravariantly on cohomology, the resulting
endomorphism algebra is the opposite algebra,
which is isomorphic to the same quaternion algebra.

The exact relations are
\[
             \kappa^2=-3,\qquad j^2=2,
                         j\kappa=-\kappa j.
                                                               \tag{A.51.13}
\]
The four operators \(1,\kappa,j,\kappa j\)
are rationally independent: the first two are
\(F\)-linear and the last two \(F\)-antilinear,
and neither pair has a nonzero relation.
This proves the actual multiplication (A.51.1).

\paragraph{The division algebra forces the required reduction at three.}
The quaternion algebra \((-3,2)\) is division
over \(\mathbf Q_3\). To check this directly,
\(\mathbf Q_3(\sqrt{-3})/\mathbf Q_3\) is
ramified. A unit norm reduces to a square in
\(\mathbf F_3^\times\), hence to one; two is
not a norm. In the quaternion presentation the
reduced norm of \(a+bj\) is
\(N(a)-2N(b)\). A nonzero zero of this norm
would make two a norm, so the reduced norm is
anisotropic and the algebra is division.
Every nonzero module over this division algebra
has \(\mathbf Q_3\)-dimension at least four.

Pass to a finite local extension over which
\(B\) is semistable and all its endomorphisms
are defined. Endomorphisms extend to the Neron
model and preserve its torus. If that torus
had positive rank \(r\le2\), its character
space tensored with \(\mathbf Q_3\) (or the
contravariant dual) would be a nonzero unital module over the
four-dimensional division algebra. This is
impossible. The toric rank is therefore zero,
so semistability is good reduction. This is
the same extension-and-specialization argument
used in A.44, now with a division algebra in
place of a full CM field; the extension facts
are those in \cite[proof of Theorem 6(a), p.~503]{fghSerreTate1968}.

On that good special fibre the rational etale
Tate space of its three-divisible group has
dimension equal to the three-rank, at most two.
It too is preserved by every endomorphism.
If nonzero it would be a unital module over
\((-3,2)_{\mathbf Q_3}\), again impossible.
Thus the special three-rank of \(B\) is zero.
The word ``unital'' matters: the identity
endomorphism restricts to the identity on any
nonzero toric or etale part, so this cannot be
an action which kills the division algebra.

The original \(E\) has good supersingular
reduction at three. The CM elliptic curve
\(E_0\) has potentially good reduction by
\cite[Theorem 6(a)]{fghSerreTate1968}; its ramified
quadratic CM field over \(\mathbf Q_3\)
cannot act on a nonzero etale three-Tate space
of dimension at most one, so its special
three-rank is zero as well. By (A.50.2) the
same potential good reduction and three-rank
zero hold for \(J_C\).

\paragraph{All cusp differences and the full integral fixed subgroup.}
For each \(T\in E[3]\), let \(P_T\) be
its unique preimage on \(C\). If \(T,U\)
are two such points, \(3(T-U)\) is principal
on \(E\). Total ramification of \(q\)
gives
\[
             q^*\bigl(3(T-U)\bigr)
                           =9(P_T-P_U).             \tag{A.51.14}
\]
Thus every cusp difference is killed by nine,
not merely the sum of the eight finite cusps.
After good reduction of \(J_C\) at three,
all these differences specialize to zero,
since its special three-rank is zero.

The stable reduction of \(C\) is of compact
type, because its Jacobian is good. Separate
all cusp sections in a semistable pointed model.
The Abel-map criterion used in A.42 shows that
any two of these marks are joined entirely
through rational components, including their
marked endpoint components: the only connected
fibres of the special Abel map are the rational
subtrees. Equivalently one may apply directly
the p-primary cusp-difference sufficient
criterion of A.42. There are nine distinct
cusps, so every selected cusp has another one
to which this criterion applies.
A.42 then proves the full fixed inertia assertion
at every coefficient prime, for every indicated
packet and open derived subgroup. This final
step is the proved group theorem, including its
integral cusp-power subtraction; no rational
Malcev fixed element is silently lifted.
\end{proof}
